\documentclass[11pt]{article}
\usepackage[letterpaper,margin=.84in]{geometry}
\usepackage{amsmath,amssymb,booktabs,microtype,hyperref,xcolor,fancyhdr}
\ifdefined\pdfmapfile\pdfmapfile{+cm.map}\pdfmapfile{+cmextra.map}\pdfmapfile{+symbols.map}\fi
\definecolor{ink}{HTML}{193D53}\hypersetup{colorlinks=true,urlcolor=ink,linkcolor=ink,pdftitle={Prime interactions in irreducible covering sets},pdfauthor={Codex, OpenAI}}
\pagestyle{fancy}\fancyhf{}\lhead{Erd\H os problem 1189}\rhead{Research manuscript for review}\cfoot{\thepage}
\setlength{\headheight}{14pt}\setlength{\parindent}{0pt}\setlength{\parskip}{5pt}
\newcommand{\head}[1]{\par\vspace{4pt}{\large\bfseries\color{ink}#1}\par}
\newcommand{\eps}{\varepsilon}
\newcommand{\lp}{P^+}
\begin{document}
{\LARGE\bfseries\color{ink} Prime interactions in irreducible covering sets}\par
{\large A necessary dense subgraph for a gain in reciprocal sum}\par
{Codex (OpenAI), at Jordan Pollard's request.\hfill 27 November 2026}

Let $A$ be a set of $k$ distinct integers greater than one. It covers if one residue per modulus can cover $\mathbb Z$. It is irreducible if every proper subset fails for every choice of residues. Write $M=\max A$, $L=\operatorname{lcm}(A)$, $S=\sum_{n\in A}1/n$. The original questions ask for the least $M$ and greatest $S$ at each size.

Assign each $n$ a divisor $b(n)>1$, with at most $b$ moduli assigned the same $b$. Such labels always exist, for example $b(n)=n$. On the primes join $p,q$ if they both divide at least one label. Its degeneracy is the largest minimum degree of any nonempty subgraph, with isolated vertices allowed. Let $D$ be the least possible degeneracy over all assignments. Let $D_{<}$ instead minimize the maximum number of \emph{larger} neighbors of a prime. Both can be zero, as prime powers create no edges.

\head{Two bounds for unrestricted irreducible sets}
For $V=\log(2+D)+2\log\log(3k)+2$, one has
\[
 S\le \tfrac12\log k+
 O\!\left(V\left[1+\log^2\!\left(2+\frac{\log k}{V}\right)\right]\right).\tag{1}
\]
In particular, any fixed gain $S\ge(1/2+\eps)\log k$ forces $D\ge k^{c(\eps)}$ eventually. \emph{Every} feasible choice of labels then contains a subgraph in which every participating prime has polynomially many neighbors. Stars of arbitrary size and all graphs of subpower degeneracy preserve the sharp leading coefficient $1/2$ for classes containing the earlier every-size construction.

For each fixed $C>0$, among irreducible sets with
\[
 \log(2+D_{<})\le C\log k/\log\log k
\]
the minimum satisfies $M_{\min}(k)\sim k\log k$. This also holds for every intermediate class containing the earlier constructions and satisfying this bound. In particular it holds for iterated substitutions on disjoint bounded-size prime sets, once the zero-defect constructions are included (page~4). The restriction permits growing sets of interacting primes.

Both obstructions depend on optimizing the choice of labels. At each size from 26 through 30 the previously exhibited extremizer admits only-prime-power labels. Page~4 gives irreducible examples where this fails for nearly all rows and one edge still suffices. A positive gain also forces complete products of two prime sets and a common cofactor set, all three of size at least $c_\eps\sqrt{\log k}$ (pages 5--6). For a fixed set of hub primes we compute both the exact leading counting constant and the reciprocal expansion through a bounded remainder in the relaxation with coverage (pages 7--12). The unrestricted optimal coefficient and minimum remain open.

\newpage
\head{1. The reciprocal bound, including arbitrary orientations}
A graph of degeneracy $D$ has an orientation with at most $D$ outgoing edges at each vertex, by successively deleting a vertex of minimum degree. Fix an optimal labeling and this orientation. Call $n$ exceptional if $\lp(n)\le\sqrt n$. All other $n$ equal $pd$ with $1\le d<p$ and $p\mid L$. They contribute at most $\sum_{p\mid L}H_{p-1}/p$, where $H_i=\sum_{j\le i}1/j$. The Zn\'am--Simpson bound gives $\sum_{p\mid L}(p-1)\le k-1$. Up to $Y=\sqrt{k\log k}$ the prime sum is at most $\log Y+O(\log\log(3Y))$, by the classical Mertens estimates. Beyond $Y$ it is $O(k\log Y/Y^2)=O(1)$, comparing each summand to $p-1$. This gives $\frac12\log k+O(\log\log(3k))$.

For $2\le t\le k$ let $E(t)$ count exceptional rows up to $t$. Put $Z=(D+2)\log^2(3k)$ and $x=\log t$. Rows whose label has two prime factors exceeding $Z$, counted with multiplicity, occupy a fraction $O(1/\log(3k))$ of $[1,t]$. Indeed repeated primes cost $\sum_{p>Z}p^{-2}\ll1/Z$. For distinct factors, orient each edge and sum $1/(pq)$ from the tail $p$; at most $D$ heads contribute, each $1/q\le1/Z$. The remaining sum is at most $(D/Z)\sum_{p\le t}1/p$. This uses no ordering of the endpoints.

Otherwise write the label as $e$ or $ep$, with all primes of $e$ at most $Z$ and $p>Z$ prime. Let $n_Z=\prod_{q\le Z}q^{v_q(n)}$. For $s=1/(2\log(2Z))$ the uniform estimate
\[
 t^{-1}\sum_{n\le t} n_Z^s\le
 \prod_{q\le Z}\left(1+\frac{q^s-1}{q-q^s}\right)\ll1 \tag{2}
\]
follows by expanding positive divisor increments $q^{js}-q^{(j-1)s}$ and bounding the frequency of multiples by $1/q^j$. For the product use $q^s-1\ll s\log q$, $q-q^s\gg q$ and $\sum_{q\le Z}\log q/q\ll\log(2Z)$. Thus $e>t^\eta$ costs proportion $O(\exp(-\eta x/(2\log(2Z))))$.

With $e\le t^\eta$ and $0<\eta<1/10$, labels $e$ or $ep$ for $p\le t^{1/2-2\eta}$ contribute $O(t^{2\eta}+t^{1-2\eta})$ rows by their capacities and summing all possible $e$ and $p$. The remaining $p$ divides the exceptional $n$, hence $t^{1/2-2\eta}<p\le\sqrt t$. Its multiples cost proportion $O(\eta+1/x)$, by partial summation of $\pi(v)\ll v/\log v$ on that interval.

Take $\eta=8V\log(2+x/V)/x$. When $\eta\ge1/10$ the trivial bound $E(t)\le t$ costs $O(V)$ in the integral, since then $x=O(V)$. Else the preceding bounds give
\[
 E(t)/t\ll (\log(3k))^{-1}+(1+x/V)^{-2}
                  + V\log(2+x/V)/x .
\]
Partial summation to $k$ and integration in $x$ give the error in (1). The $k$ rows above cutoff $k$ contribute at most 1 in total. Finally $z(1+\log^2(2+1/z))\to0$ as $z\downarrow0$, proving the polynomial-degree necessity.

\newpage
\head{2. The largest modulus with ordered interactions}
Let $C(t)$ count rows up to $t$ with a composite label, under a labeling attaining $d=D_{<}$. For every fixed $a>0$,
\[
 C(t)/t\ll_a t^{-1/3}+e^{-\sqrt{\log(2t)}}+
 \exp\!\left(\frac{-a\log t}{\log(2+d)+\sqrt{\log(2t)}}\right). \tag{3}
\]
Set $B=(d+2)\exp(\sqrt{\log(2t)})$. If all factors of a composite label exceed $B$, it contains $p^2$ or $pq$ with $B<p<q$ and $q$ a larger neighbor of $p$. Their multiples cost proportion $\le(d+1)\sum_{n>B}n^{-2}\ll(d+1)/B$. Otherwise its least prime is at most $B$ and every one of its prime factors belongs to a set $T$ of at most $(d+1)\lfloor B\rfloor$ primes: the small primes and their larger neighbors. Labels of size $\le t^{1/3}$ contribute $O(t^{2/3})$ rows by capacity. The rest require the $T$-part of their modulus to exceed $t^{1/3}$.

To bound that event let $H=(2+|T|)^2$ and $s=\min(1/2,A/\log H)$ for any fixed $A>0$. The proof of (2) applies to $T$ with constant depending only on $A$. For primes $q\le H$ use $q^s-1\ll_A s\log q$; those above $H$ add $O(|T|H^{-1/2})$ to the logarithm. Markov's inequality gives the fraction $O_A(t^{-s/3})$. Since $\log H\ll\log(2+d)+\sqrt{\log(2t)}$, taking $A$ sufficiently large in terms of $a$ proves (3); bounded $t$ and the $s=1/2$ range are absorbed into the constant.

For the remaining prime labels at $t=M$, each $p$ supports at most $p$ rows. Each active $p\mid L$ is $\le\sqrt M$: within a fibre over $L/p$, a private point of a highest-$p$-power row implies at least $p$ such rows, all distinct multiples of $p^{v_p(L)}$. Hence the number of prime-labeled rows is at most $\sum_{p\le\sqrt M}p\sim M/\log M$. If $M>2k\log k$ the desired lower bound already holds. Otherwise $\log M\sim\log k$, and (3), with a sufficiently large fixed $a$ depending on $C$, gives $C(M)=o(M/\log M)$. The zero-defect constructions in the companion paper attain both leading bounds at all sufficiently large sizes using the labels $p^j$ alone.

\head{Why the ordering matters for this proof}
For a real cap $X$ take odd primes $p\le\sqrt X$, all $pd$ for $1\le d<p$, and $2pe$ for $p/2\le e<p$ with $2pe\le X$. Include the fixed cover $2,3,4,6,12$, omitting duplicates. The displayed families are distinct: $p$ is the largest prime and their cofactors are respectively below and above $p$. Label them by $p$ and $2p$; use $2$ and $4$ for the two additional rows when $X$ is large. This is a star. Partial summation gives $(1+\tfrac12\log2+o(1))X/\log X$ rows. Indeed the extras below $\sqrt{X/2}$ contribute $1/4$ in these units, and those above contribute $\int_{1/\sqrt2}^1(r^{-1}-r)\,dr$. The fixed cover is a proper subset. Thus discarding irreducibility while retaining sparse unordered interactions and coverage cannot give the stated cardinality estimate. Whether the star suffices for a genuine irreducible improvement is open here.

\newpage
\head{3. Irreducible substitutions and a one-edge family}
Distinguish one indexed row $d$ in an irreducible multiset on an old period $L_0$. Let $b_i>1$ be an irreducible indexed multiset on a period coprime to $L_0$. Replace the distinguished row by $b_i e_i$, with $e_i\mid d$ and $\operatorname{lcm}_i e_i=d$. The resulting multiset is irreducible under arbitrary reassignment. To obtain pairwise distinct moduli, check distinctness at the end.

Indeed in any assignment the old rows leave holes. Above each hole the active new rows must cover the entire new factor. By irreducibility this requires every index $i$. All old holes therefore lie in the intersection of the exterior classes modulo $e_i$, which, when nonempty, is one class modulo $d$. Conversely align all exterior residues with a chosen $d$-class and use an internal cover; CRT supplies the new residue for each row. A deletion outside the packet contracts to the old deletion. After an inside deletion the other new rows cannot fill the fibre over any old hole. This argument also works with the initial virtual row $d=1$. The accompanying Lean development proves this substitution for actual integer congruences, using the lcm divisibility property and integer Bezout coefficients for the coprime periods. Its conclusion quantifies over arbitrary subfamilies and all new assignments of residues.

At each substitution use the internal modulus as label on each surviving new row. An internal modulus $b$ occurs at most $b$ times, since $b$ copies can already cover. All primes in that step are fresh and the periods persist after contraction, so labels from different steps cannot coincide. The graph is a union of graphs on the prime sets of the separate steps. This permits arbitrary local irreducible multisets within each step.

For an explicit family let $p,q$ be distinct primes in $[z,2z]$ with $z$ growing. Set $t=\lceil\log_2(pq)\rceil$ and use the first $t$ primes $r_1,\ldots,r_t$ as seed, all less than $p,q$ eventually. Start with one virtual row $d=1$. At prime $r_i$ replace it by $r_i$ copies of the internal modulus $r_i$, dressing the first $r_i-1$ with distinct divisors of $d$ and the last with $d$. Continue on the latter row, now $r_i d$. There are enough divisors since $2^{i-1}\ge r_i-1$ by Bertrand's postulate and induction. Interim coincidences with the distinguished row are permitted. At the end use $pq$ copies of $pq$, with distinct seed divisors including the full seed period. All final rows are distinct by their newest prime and their cofactor. The resulting size is $k=pq+\sum_i(r_i-1)$. Every prime-power label is just a prime here, so the total such capacity is at most $p+q+\sum_i r_i=O(\sqrt k)$. Yet one edge, $pq$, suffices for all $k$ rows. Thus the proportion that cannot receive prime-power labels can tend to one even for irreducible sets of degeneracy exactly one.

For a smaller instance take seed $2^3 3^2\cdot11\cdot13$ and packet $5\cdot7$. A prime-power stage uses $r-1$ copies of each $r^j$ plus a distinguished copy of the final power. Its reciprocal mass is exactly one and first-nonzero digits give coverage, hence it is internally irreducible. The constructor gives 64 distinct moduli of maximum 13,860 and period 360,360. Exact matching places 62 with prime powers and all 64 with the additional edge $5,7$. Lean checks size, distinctness, maximum, and full irreducibility of the explicit set: for every predicate on its moduli, the retained moduli admit some cover of all integers exactly when every modulus is retained.

{\small\textbf{Sources and status.} \href{https://www.erdosproblems.com/1189}{Erd\H os 1189}; \href{https://doi.org/10.13140/RG.2.2.16138.63684}{E. \"Ozt\"urk, 2026} (prior upper coefficient $3/2-\log 2$); the every-size lower construction and period-budget proof in \href{https://github.com/jpollard-cs/open-math-agents/pull/1}{the companion manuscripts}. The analytic results are written proofs for review. The finite checks and kernel proof have the narrower scopes stated above.}
\newpage
\head{4. Complete products of primes and cofactors}
There is a concrete consequence for the moduli themselves. Fix $\eps>0$. If an irreducible $k$-set has $S\ge(1/2+\eps)\log k$, then for all sufficiently large $k$ it contains distinct products
\[
 \{pqd:p\in U,\ q\in W,\ d\in B\},\qquad
 |U|,|W|,|B|\ge c_\eps\sqrt{\log k}.
\]
Here $U,W$ consist entirely of primes, each in its own dyadic interval; $B$ consists of positive integers. Every product $n$ is at most $k$ and satisfies $\lp(n)\le\sqrt n$. All three factors are $\ge k^{c_\eps}$ and $q/p\ge k^{c_\eps}$. For \emph{every} feasible labeling the box can be chosen with $pq\mid b(n)$ throughout. There are such boxes in $\gg_\eps(\log k)^3$ distinct triples of dyadic scales for $n,p,q$, using disjoint sets of moduli. Constants in this section may depend on $\eps$.

\emph{Finding enough rows.} The nonexceptional rows contribute at most $\frac12\log k+O(\log\log k)$ by page 2. Discard $n>k$ at cost at most 1. Put $Z=k^\delta$, where $\delta>0$ is a sufficiently small constant in terms of $\eps$. Among the exceptional rows, those whose label has at most one prime factor above $Z$ counted with multiplicity cost at most
\[
 O\!\left((1+\log(2Z))\,[1+\log^2(2+\log k/(1+\log(2Z)))]\right).
\]
This is exactly page 2's $e,ep$ calculation with $V=1+\log(2Z)$; the oriented-edge term is absent. Rows with a squared prime $p>Z$ dividing the label contribute $\ll(1+\log k)\sum_{p>Z}p^{-2}=o(1)$. Deleting also $n<k^{\min(\eps/8,1/8)}$ costs at most $(\eps/8)\log k+O(1)$. By first choosing $\delta$ small, the remaining mass is at least $\eps\log k/2$ for large $k$. Assign each remaining row \emph{one} pair of distinct primes $Z<p<q$ dividing its label; the choice can be arbitrary. For each row $q\le\sqrt n$.

\emph{Finding dense boxes of available triples.} At least $c\log k$ dyadic bands $[T,2T)$ each have $\ge cT$ remaining rows, by their reciprocal mass and the bound of $T$ rows per band. For a fixed such $T$, sort the assigned pairs into dyadic prime bins $[P,2P)$, $[Q,2Q)$ with $P\le Q$. There are $O((\log k)^2)$ bin pairs. Each accounts for at most $C_\delta T/(\log k)^2$ rows: there are at most $C_\delta P Q/(\log k)^2$ prime pairs by Chebyshev and at most $2T/(P Q)$ multiples per pair. Bins meeting the endpoints are included; $P,Q\ge k^\delta/2$ and $P Q<2T$ whenever occupied. Consequently $\gg_\eps(\log k)^2$ bin pairs each have $\gg_\eps T/(\log k)^2$ assigned rows (discard bins below a small constant times the mean).

For each $P$ we may also discard $Q/P<k^\beta$ and $T/(P Q)<k^\beta$. These two strips together contain $O(\beta(\log k)^2+\log k)$ occupied bin pairs. For sufficiently small fixed $\beta>0$ there remain $\gg_\eps(\log k)^2$ dense pairs. The primes in each full bin and all integers in $[T/(4PQ),2T/(PQ)]$ form three vertex sets of size $\ge k^{c}$ by the prime number theorem and the width of the cofactor interval. Their product of cardinalities is at most $C_\delta T/(\log k)^2$. The assigned triples $(p,q,n/(pq))$ thus have density bounded below by a positive constant, say $\alpha$, in this Cartesian product.

\newpage
\head{Extracting a complete subproduct}
For completeness the quantitative dense-hypergraph calculation is given here. Let $X,Y,Z'$ be the three vertex sets, all of size $\ge k^c$. At least an $\alpha/2$ fraction of pairs $(y,z)$ have at least $\alpha|X|/2$ neighbors in $X$. A uniformly random set of $r$ distinct vertices of $X$ lies in each such neighborhood with probability at least $(\alpha/4)^r$, provided $r\le\alpha|X|/4$. Hence some $r$ vertices have a common bipartite graph on $Y,Z'$ of density
\[
 \lambda\ge (\alpha/2)(\alpha/4)^r .
\]
At least $\lambda|Z'|/2$ vertices of $Z'$ see at least $\lambda |Y|/2$ vertices of $Y$ in this graph. The same argument gives $r$ vertices in $Y$ with at least $(\lambda |Z'|/2)(\lambda/4)^r$ common neighbors, provided $r\le\lambda |Y|/4$. Both this condition and the requirement of at least $r$ neighbors hold for $r=\lfloor c'\sqrt{\log k}\rfloor$ with $c'>0$ sufficiently small in terms of $c,\alpha$. These vertex sets give the asserted complete product. No two triples can yield the same modulus because a modulus was assigned only one pair. The same assignment prevents repetitions between the different triples of dyadic scales.

This last extraction is the classical dense complete multipartite hypergraph argument; compare \href{https://renyi.hu/~p_erdos/1964-13.pdf}{Erd\H os, 1964}. The exponent $1/2$ is the limit of an argument from positive triple density alone: a random constant-density triple system on three parts of size $k$ has expected number of balanced boxes of side $r$ at most $k^{3r}\alpha^{r^3}$, tending to zero for $r>C_\alpha\sqrt{\log k}$. Its density concentrates. This observation gives no upper bound on the boxes of an actual irreducible cover.

\head{Where the unrestricted problem remains}
The product structure connects a proposed gain in reciprocal mass with a testable pattern in the same cover. It does not yet contradict irreducibility. In particular, an isolated displayed box can always miss a point, whatever residues are assigned: at each $p\in U$ avoid the at most $|W|\,|B|$ residues specified by the box, then combine the choices by CRT. Here $p\ge k^c$ whereas we can keep the two sides of size $O(\sqrt{\log k})$. Any argument using a box to show a row is removable must exploit rows outside it or a different feature of irreducibility. A concrete remaining question is whether the other rows can provide the required prime fibres simultaneously across a positive proportion of the three-dimensional scale ranges.

\newpage
\head{5. Exact counting cost of finitely many shared primes}
Fix a finite set of primes $H$ and put $v=\prod_{h\in H}(1-1/h)$. Consider sets of distinct integers $2\le n\le X$ with all prime divisors $\le\sqrt X$ and some feasible assignment of labels in which every label involves at most one prime outside $H$. Thus $H$ is a vertex cover of the label graph. Let $K_H(X)$ be the largest size of such a set. Then
\[
 K_H(X)\sim C_H\frac X{\log X},\qquad
 C_H=1-\log v+2\sum_{h\in H}\frac{\log h}{h-1}.\tag{4}
\]
The same constant holds on imposing coverability. It gives $M\ge(1/C_H-o(1))k\log k$ for genuine irreducible sets with this restriction. Attainment with irreducibility for nonempty $H$ remains open. At $H=\{2\}$ the exact coefficient in (4) is $1+3\log 2$, improving the elementary lower example on page 3. This quantifies the information lost by using just the label graph, distinctness, the active-prime bound, and coverability.

\emph{Upper bound.} Let $\mathcal T$ be the positive integers supported on $H$, including 1. A label $b$ occurs at most $\min(b,X/b)$ times. Pure labels in $\mathcal T$ contribute $O_H(\sqrt X(1+\log X)^{|H|})$. Labels $s p^a$ with $s\in\mathcal T$, $p\notin H$, $a\ge2$ contribute at most
\[
 O_H((1+\log X)^{|H|+1})\sum_{p\le\sqrt X}\min(\sqrt X,X/p^2)
 = O_H(X^{3/4}(1+\log X)^{|H|+1}).
\]
For fixed $p$ all remaining labels are $s p$. Separate the moduli by their full $H$-part $t$. When $t\le y=\sqrt{v X}/p$, all possible labels have $s\le y$, yielding capacity $\sum_{s\le y,\,s\in\mathcal T}sp$. For $t>y$ the moduli are multiples of $pt$ whose quotient is coprime to $H$, hence number at most $vX/(pt)+O_H(1)$ for each $t\le X$, by inclusion-exclusion. After summing the rounding errors over $p,t$ we get $O_H(\sqrt X(1+\log X)^{|H|})$. The resulting upper bound up to $o(X/\log X)$ is
\[
 \sum_{s\in\mathcal T}\sum_{p\le\sqrt X}\min(sp,vX/(sp)).
\]
For fixed $s$, split the prime sum at $\sqrt{vX}/s$. By the prime number theorem and partial summation its coefficient in units $X/\log X$ tends to
\[
 \frac vs\left(1+\log(s^2/v)\right).\tag{5}
\]
For $s\le X^{1/4}$ the normalized sum is $O_H((1+\log s)/s)$: up to the cutoff apply $\sum_{p\le z}p\ll z^2/\log(2z)$, and beyond it sum $1/p$ in dyadic intervals, all above a constant times $X^{1/4}$. For $s>X^{1/4}$ use the unnormalized bound $O(X\log\log(3X)/s)$ and $\sum_{s>X^{1/4},\,s\in\mathcal T}1/s\le X^{-1/8}\sum_{s\in\mathcal T}s^{-1/2}$. These justify summing (5). Since $v/s$ is a probability measure on $\mathcal T$ with independent geometric prime exponents of means $1/(h-1)$, its expectation of $\log s$ is $\sum_{h\in H}\log h/(h-1)$. This proves the upper coefficient in (4).

\newpage
\head{Constructions attaining the counting coefficient}
Fix $K$ and $a\in(0,1)$, take $s\in\mathcal T$, $s\le K$, and primes $a\sqrt X\le p\le\sqrt X$ outside $H$. Use up to $sp$ distinct integers $n=spe\le X$ with $e$ coprime to all of $H$ and every prime dividing $e$ strictly smaller than $p$. Assign the label $sp$. The full $H$-part identifies $s$ and the largest prime outside $H$ identifies $p$, so the choices cannot collide across rows. The number of eligible $e$ equals
\[
 vX/(sp)+O_{H,K,a}(1+p/\log p).
\]
Indeed $e\le a^{-2}p/s$, and the rejected integers containing a prime $\ge p$ have the form $jb$ with $j\le a^{-2}$ and $b$ prime of size $O_a(p)$. The total error in the count is $o(X/\log X)$. As $a\downarrow0$ and then $K\to\infty$, these finite truncations attain (4). For large $X$ the constructed moduli all exceed 12. Add $2,3,4,6,12$, labeled respectively by $2,3,4,3,4$, to include a fixed cover; no established label capacity is exceeded. This supplies the matching lower bound with coverage.

For nonempty $H$ the construction can also obey the necessary Zn\'am--Simpson period budget. In the fixed truncation, primes occur only through $\sqrt X$, with first-power budget at most $(1+o(1))X/\log X$. Repeated powers outside $H$ come from $e=O_a(\sqrt X)$ and their extra budget totals $O_a(\sqrt X\log X)$, by summing over primes $O_a(X^{1/4})$; powers in $H$ are bounded by $K$. Since $C_H>1$ we may take truncations of coefficient strictly above one and arbitrarily close to $C_H$. For $H=\varnothing$ the earlier irreducible constructions already attain coefficient one at leading order (and satisfy the budget).

\head{A remaining question with one central prime}
For the single central prime 2, (4) locates the unrestricted irreducibility issue between the sharp arithmetic coefficient $1+3\log 2$ and the achieved irreducible coefficient one, in units $X/\log X$ for counts. Can an irreducible family use the same binary coordinate at many odd primes to raise the coefficient above one? Projecting onto the two parity classes exposes the difficulty. Even rows can be reassigned between the classes, whereas every odd row contributes to both. A proposed construction must survive arbitrary redistribution after every deletion. The examples attaining (4) have a five-row proper cover, so they give no irreducible improvement.

There is an exact finite identity behind the cutoff. For a fixed exterior prime $p$, allow all $pm$ with $m\le N$, disregarding other candidate primes. If $u_s$ counts $m$ with exact $H$-part $s$, the placement maximum is $\sum_s\min(ps,u_s)$. The lower bound uses the exact part as label cofactor. For the upper bound, $u_s$ counts integers up to $N/s$ coprime to $H$, hence decreases with $s$, whereas $ps$ increases. Up to their crossing all labels available to a row of exact part $s$ have cofactor at most $s$, so use the total capacity on that initial segment. On its complement count the rows. This includes the cases of an empty initial segment or an exact tie.

The included program independently compares direct bipartite placement with the cutoff calculation for small caps. It also counts the labeled construction at larger cutoffs. These are finite checks of the optimization mechanism. The limit (4) follows from the written argument.

\newpage
\head{6. The same constant in the second reciprocal term}
Fix $H$ as on page 7. Among $k$ distinct moduli with a label graph covered by $H$, satisfying the period budget $\sum_{q^a\parallel L}a(q-1)\le k-1$ and admitting a cover, let $R_H^*(k)$ denote the largest reciprocal sum. For nonempty $H$, and also for empty $H$ by the earlier zero-defect construction, we obtain
\[
 \lim_{k\to\infty}\frac{R_H^*(k)-\tfrac12\log k}{\log\log k}
    =\tfrac12+\gamma+C_H-1.\tag{6}
\]
Both bounds have remainder $O_H(1)$ after the two displayed orders, as shown on page 11. It applies to genuine irreducible sets. The construction for nonempty $H$ establishes the lower bound for the relaxation with coverage. In particular for one binary hub the possible additional second-order coefficient is bounded by $3\log2$ under irreducibility, and equals $3\log2$ in this relaxation.

Pure $H$-labels contribute at most $\sum_{s\in\mathcal T}H_s/s=O_H(1)$. Labels $sp^a$ with $a\ge2$, $p\notin H$, contribute at most
$\sum_{s,p,a\ge2} H_{sp^a}/(sp^a)=O_H(1)$, where only the exponent is restricted to be at least two. This follows from convergence of $\sum_{s\in\mathcal T}(1+\log s)/s$ and $\sum_p(1+\log p)/p^2$.

For labels $sp$ with fixed $p\notin H$, the finite rank identity on page 8 is attained at \emph{every} cutoff by one infinite set: for each exact $H$-part $s$ select the $sp$ smallest $m$ with that exact part and take $n=pm$. This dominates the cumulative counts for any assignment, so partial summation bounds its reciprocal mass by that of this set, regardless of which exterior prime ultimately receives a modulus. Let $A_H(j)$ be the reciprocal sum of the first $j$ integers coprime to all of $H$. Inclusion-exclusion gives
\[
 A_H(j)=v[\log(j/v)+\gamma+\sum_{h\in H}\log h/(h-1)]+O_H(1/j).
\]
Indeed their counting function through $u$ is $vu+O_H(1)$ and their harmonic sum is $v(\log u+\gamma)-\sum_{d\mid\prod H}\mu(d)\log d/d+O_H(1/u)$. The divisor sum equals $-v\sum_{h\in H}\log h/(h-1)$ by differentiating the finite Euler product. The local optimal mass therefore equals
\[
 \frac1p\sum_{s\in\mathcal T}\frac{A_H(sp)}s
   =\frac{\log p+\gamma+C_H-1}{p}+O_H(p^{-2}).
\]
Sum through $Y=\sqrt{k\log k}$ and bound larger active primes by the period budget as on page 2. Mertens' estimates give the upper bound in (6), including the stated bounded remainder.

\newpage
\head{Exact cardinality and the matching lower construction}
Take $H\ne\varnothing$, fix $J\ge1$ and put $B=\prod_{h\in H}h^J$. The condition $\gcd(m,B)=s$ for $s\mid B$ has density $w_s=\varphi(B/s)/B$ with a bounded counting error depending only on $B$. For each prime $p>\max(H\cup\{13\})$, and each $s\mid B$, select the $sp$ smallest integers $m$ with this gcd, then discard those with a prime factor $\ge p$. Put $pm$ into the family with label $sp$. Before deletions this includes every $m\le p$, since no gcd class has more than $p$ entries up to $p$. Its endpoint for gcd $s$ is $(s/w_s)p+O_B(1)$. Partial summation on $p<m\le(s/w_s)p+O_B(1)$ gives additional harmonic mass $w_s\log(s/w_s)+O_B(1/p)$ in the $m$ variable. The $O_B(p/\log p)$ rejected $m$, all at least $p$, cost $O_B(1/(p\log p))$ in the original moduli. Thus the row block at $p$ contributes
\[
 \frac{\log p+\gamma+g_B}{p}+O_B(1/(p\log p)),\qquad
 g_B=\sum_{s\mid B}w_s\log(s/w_s).
\]
No modulus occurs twice, since $p$ is its largest prime, of exponent one. Each label has the required capacity. Include the fixed five-row cover with the labels given on page 8. Process $p$ in increasing order and take any prefix of the final block to realize each sufficiently large cardinality. Through a full block $x$, the size is $(\sum_{s\mid B}s+o_B(1))x^2/(2\log x)$. The terminal partial block costs $O_B(x)$ rows and its omission loses $O_B(\log x/x)$ harmonic mass, using the prime number theorem to compare successive primes. Since $\sum_{s\mid B}s>1$, the first-power budget $\le(1+o(1))x^2/(2\log x)$ leaves room for all repeated powers, whose extra cost is $O_B(x\log x)$ as on page 8. Consequently at every sufficiently large exact $k$,
\[
 S=\tfrac12\log k+(\tfrac12+\gamma+g_B)\log\log k+O_B(1).
\]
The bounded error follows by summing the displayed block estimate, since $\sum_p1/(p\log p)$ converges, and by taking logarithms in the size asymptotic.

The distribution $w_s$ is that of independent geometric prime exponents, each stopped at $J$. In fact
\[
 g_B=\sum_{h\in H}(1-h^{-J})\left[-\log(1-1/h)+2\log h/(h-1)\right].
\]
Given that an exponent reaches $J$, its remaining length has the original geometric law. Truncating it therefore loses the fraction $h^{-J}$ of both its mean and its entropy (by the entropy chain rule). Therefore $g_B$ is the expected log of the resulting integer plus the entropy of its distribution (natural logarithms). As $J\to\infty$ the expectation and entropy converge to those of the independent untruncated geometric variables; their geometric tails times $1+j$ are summable. This gives $g_B\to \sum_{s\in\mathcal T}(v/s)\log(s^2/v)=C_H-1$. First take $k\to\infty$ for fixed $J$, then $J\to\infty$. This proves (6). Page 11 refines the choice of cutoffs to control the remainder. For $H=\varnothing$ the accompanying zero-defect paper gives the asserted lower expansion.

\newpage
\head{Uniformly bounded remainder at every exact size}
The two-cutoff construction below strengthens the lower estimate in (6) to
\[
 R_H^*(k)=\tfrac12\log k+
       (\tfrac12+\gamma+C_H-1)\log\log k+O_H(1).\tag{7}
\]
The extra rows needed for the reciprocal term can be spent on the smaller exterior primes without changing the final cardinality scale.

For a prescribed sufficiently large $k$, put $y=k^{1/8}$. For each sufficiently large prime $p\le y$ outside $H$, use the preceding gcd construction with $J(p)=\lfloor c_H\log\log p\rfloor$, choosing the positive constant $c_H$ so that $B(p)=\prod_{h\in H}h^{J(p)}\le(\log p)^{1/16}$. Omit finitely many starting primes, keeping the fixed cover. For all $p>y$ use one fixed $B_0=\prod_{h\in H}h$, continuing in increasing order until size $k$ is reached.

Here are uniform estimates for the small primes. A gcd class mod $B$ has counting error at most $B$ and density $w_s\ge1/B$. Its endpoint is $sp/w_s+O(B/w_s)$ and is $O(B^2p)$. Its excess $m$-harmonic mass beyond $p$ differs from $w_s\log(s/w_s)$ by $O(B/p)$ by partial summation (the endpoint error is multiplied by $w_s$). Across all $s\mid B$ this costs $O(B^2/p)$. Every rejected $m$ is at least $p$ and equals $jr$ with $j\ll B^2$, $r\ge p$ prime. By Chebyshev the number is $O(B^2p(1+\log(2B))/\log p)$, summing $1/j$ for these $j$. Thus the errors in the original reciprocal sum, after division by $p$, are bounded in total, since
\[
 \sum_p\left(\frac{B(p)^2}{p^2}+
       \frac{B(p)^2(1+\log(2B(p)))}{p\log p}\right)<\infty .
\]
Moreover $0\le C_H-1-g_{B(p)}\ll_H(1+J(p))2^{-J(p)}$ by the geometric tail formula for the stopped exponents and their entropies. The sum of these errors weighted by $1/p$ also converges, by partial summation and $J(p)=c_H\log\log p+O(1)$.

At most $O(y^2(1+\log y)^2)=o(k)$ rows have been spent through $y$; indeed each prime contributes at most $p\sum_{s\mid B(p)}s\le p B(p)^2$. The largest exterior prime $x$ reached is again asymptotic to a positive $H$-dependent constant times $\sqrt{k\log k}$ by the fixed-block size estimate. The partial final block and the budget are handled as on page 10, since $B(p)^2 p=o(x)$ throughout the small-prime range. In the large-prime range the missing local second-order mass is bounded by $O_H(\sum_{y<p\le x}1/p)=O_H(1)$ because $\log x/\log y$ is bounded. Summing the block estimates therefore yields (7) at the prescribed size. For the empty set $H$ use the prior construction directly. An explicit constant with an $o(1)$ remainder and the second-order term for genuine irreducible covers with nonempty $H$ are separate questions.

\newpage
\head{7. Allowing each cover to choose its own hubs}
Both optimizations remain sharp when only the number of central primes is prescribed. Let $r$ be fixed and $H_r$ the first $r$ primes. In (4) and (7), allow any labeling whose graph has vertex cover number at most $r$, with the hubs chosen afresh at each size or cap. Then the exact leading counting coefficient is $C_{H_r}$, and the reciprocal expansion is
\[
 \tfrac12\log k+(\tfrac12+\gamma+C_{H_r}-1)\log\log k+O_r(1)
\]
under the conditions of (7). The lower constructions with $H=H_r$ apply directly. In particular a single hub chosen without restriction has the binary-hub constants. These are the exact relaxation optima. They are upper bounds for the corresponding irreducible optima.

Here is the uniformity needed in taking this maximum. For $|H|\le r$ one has $v\ge\prod_{h\in H_r}(1-1/h)>0$, the inclusion-exclusion errors in counts are $O(2^r)$, and $|\mathcal T\cap[1,X]|\le(1+\log_2 X)^r$. Moreover
\[
 \sum_{s\in\mathcal T}s^{-1/2}\le (1-2^{-1/2})^{-r}.
\]
Consequently the normalized tail for $s>K$ in the proof of (4) tends to zero as $K\to\infty$, uniformly in $H$ and large $X$: for $s\le X^{1/4}$ use $(1+\log s)/s\ll K^{-1/4}s^{-1/2}$ for $s>K$, and beyond $X^{1/4}$ use the original power saving. For each fixed $s\le K$, the prime number theorem in (5) is uniform for $v$ in the indicated compact interval: all cutoff endpoints are comparable to $\sqrt X$ with constants depending only on $K,r$. Omitting exterior primes that happen to be in $H$ can only decrease the upper bound. Thus the upper bound $K_H(X)\le [C_H+o_r(1)]X/\log X$ is uniform.

For the reciprocal upper bound, inclusion-exclusion for the harmonic sum through $u$ has error $O_r(1/u)$ even when some divisors of $\prod H$ exceed $u$. For such a divisor $d$, the formal main term $(\log(u/d)+\gamma)/d$ is $O(1/u)$, since $(1+\log t)/t$ is bounded for $t\ge1$. Counting inversion is uniform since $v$ is bounded below. Thus the remainder in the expansion for $A_H(j)$ is $O_r(1/j)$. All the convergent prime-power and $\mathcal T$ tails on page 9 are likewise $O_r(1)$ by the half-power sum above. This proves a uniform $O_r(1)$ in the upper half of (7).

Finally $f(h)=-\log(1-1/h)+2\log h/(h-1)$ is strictly decreasing for $h>1$. The second term decreases because $1-1/h<\log h$. Hence the maximum $C_H$ occurs for precisely $H_r$ among sets of size at most $r$. For $r\ge1$ every other choice loses at least $f(p_r)-f(p_{r+1})>0$. It follows that a sequence within $o(X/\log X)$ of the optimal count, or within $o(\log\log k)$ of the optimal reciprocal sum, must use precisely the first $r$ hub primes in every witnessing set of at most $r$ hubs, eventually. The stability assertion concerns the stated relaxations.

\end{document}
